\section{What can and cannot be concluded}\label{sec:conclusions}

The raw difference between group time and solo flights is real, and its long-lag part survives the
tests (Figure~\ref{fig:estimates}), but where the route can be matched, most of it goes with the
route. With the pilot fixed, or with the day, site and hour fixed, a group progresses straighter over
a quarter of an hour (Section~\ref{sec:other}); reweighting the solo flights to the zones and hours
where groups fly lowers the long-lag gap by at most a fifth (Section~\ref{sec:zones}); on the half of
the group time that can be matched, fixing the route lowers it by about \RouteExplained\%
(Section~\ref{sec:route}); on the same flight the straighter progress starts and stops with the group
(Section~\ref{sec:prepost}). At
100~s, within one glide or climb, the progress is no straighter: the difference appears only at lags
of one to two climb--glide cycles. Faster and flatter glides pass the contrasts and the
before-and-after test; the climbs pass the contrasts but change only when the group forms. Two pilots
carry most of it (Section~\ref{sec:inside}).

\begin{figure}
\centering
\includegraphics[width=\textwidth]{report_estimates.pdf}
\caption{Every estimate of the long-lag difference, from the raw one to the four tests, with 5--95\%
bands. The route rows refer to the matched pieces only. Values, with the drifts of the
before-and-after test, in Table~\ref{tab:estimates}.}
\label{fig:estimates}
\end{figure}

\medskip
\noindent
\begin{minipage}[t]{0.49\textwidth}
\textbf{Supported}
\begin{itemize}
  \item the progress over a quarter of an hour is straighter in a group, with the pilot fixed
        (\DHlocThousandFlight) and with the day, site and hour fixed (\DHlocThousandGroupSdh);
  \item where the route can be matched, fixing it lowers this advantage by about \RouteExplained\%; the
        rest goes with flying together;
  \item no increase of $\Hloc(100~\mathrm{s})$, within one glide or climb;
  \item on the same flight the straighter progress starts and stops with the group, and before and
        after are alike;
  \item faster and flatter glides, within strata and on the same flight both ways;
  \item more time climbing, more searches that end in a climb and a higher climb rate within strata
        (pooled, $\G$ climbed slower than $\A$); on the same flight these rise when the group forms but
        do not fall when it breaks;
  \item selection explains little of the straighter progress, over a third of the searches that end in
        a climb and over half of the faster glides;
  \item the leader is a little faster and ends searches in a climb as often as the follower.
\end{itemize}
\end{minipage}\hfill
\begin{minipage}[t]{0.49\textwidth}
\textbf{Not supported}
\begin{itemize}
  \item that the group \emph{causes} the effect: the tests fix pilot, day, site, hour, zone and the
        drift of the flight and the route flown, but not the conditions of the day that bring pilots
        together (a line or a thermal working that day) nor the coordination needed to stay in contact;
  \item that larger groups progress straighter than pairs at the same site-day and hour, up to six
        pilots.
\end{itemize}
\medskip
\textbf{Caveats}
\begin{itemize}
  \item the glide statistics are partly fixed by definition, since headings must be aligned in glide
        for a contact (only between flights with phases, logged every second);
  \item groups of 7 or more may be competition gaggles, which the archive does not flag;
  \item the route test compares different days, and pilots who never meet anyone with pilots in
        groups;
  \item the leader is assigned from $\ell_{\rm med}$, which is uncertain in thermals (note,
        Appendix~B).
\end{itemize}
\end{minipage}

\paragraph{Open.} The phases exist on 1~Hz flights only: the same analyses on the tracks of
xc-dynamics, with phases on all loggers, and on hang gliders, which have few group flights
(Appendix~\ref{app:curves}), remain to be done. The reproduction of the report is in
Appendix~\ref{app:reproduce}, and the checks of the viewer and of the analyses in
Appendix~\ref{app:audit}.
